\section{Discussion}
\label{sec:discussion}

\begin{stub}{Haley first draft, then both; about 1.5 columns}
\emph{What was found}, in three beats: the four profiles are there and reproduce; they reproduce the way a cut continuum does; severity follows the developmental axis, not class. Two dials, not four boxes (\texttt{docs/team-notes.md} has the short version).

\emph{Against the literature.} The contrasts clinic studies report (developmental, psychiatric, high burden, few conditions) appear in a national sample, as quadrants of two axes. Litman et al.'s 37/19/34/10 against our 18/33/34/16 \citep{litman2025}; their ``broadly affected'' class and Vargason's high-rate group \citep{vargason2019} match our unstable $K=5$ class. \citet{zhang2025}'s two genetic factors as the counterpart of our two axes, including the age pattern (psychiatric position rises with age). \citet{kim2019} and \citet{wittkopf2023} reached dimensional or hybrid answers for core traits; ours is the first such comparison for co-occurring conditions that we found. The latent density result (lumpier than a bell curve, half as lumpy as four groups) fits the hybrid reading.

\emph{What replication can show.} The calibration is the methodological point: held-out-year, forward and age-band replication pass for a partition of a smooth continuum, so passing them is necessary but not sufficient evidence of subtypes. Other studies relying on these checks could run the same simulation.

\emph{Why distance methods fail on sparse binary items}: chaining in average linkage, near-tied $k$-means optima; what that means for cross-method replication as a criterion.

\emph{Limitations.} Parent report of a provider's diagnosis, so access to care and diagnostic practice are in the flags; one observation per child, so no temporal or predictive validation; entropy 0.72; 8.4\% with blank answers (refit says it does not matter); ``currently has'' wording; the $K$ rule was set after seeing results and the 2025 test is confirmatory; the proposal's hierarchical method failed and was replaced; severity and flags share an informant.

\emph{Next.} A factor mixture model \citep{lubke2005}, which allows classes and continuous variation within them; three-step corrections for relating classes to outcomes \citep{vermunt2010}.

\emph{For a clinician}: describe a child by where they sit on two dials rather than which of four boxes; severity ratings track the developmental dial.
\end{stub}
